\section{面试能力框架}
\label{sec:interview}

前十节是知识；本节把它们组织成面试中可以被观察到的能力。MLSys 面试通常同时考两件事：\emph{工程能力}——能否把一个系统问题拆成可测量、可估算的部分，
并按正确顺序定位瓶颈；\emph{研究能力}——能否读懂一篇论文的真实贡献与前提，提出能区分假设的实验。二者共享同一个底座：
式~\eqref{eq:three} 的三项分解与第~\ref{sec:perf} 节的 roofline。

\subsection{按岗位的能力矩阵}

表~\ref{tab:roles} 按五类常见岗位列出工程与研究两条线上的核心要求、必须当场算出的数，以及面试官最常深挖的方向；
它与图~\ref{fig:map} 的课程分层一一对应：推理岗主要对应第~\ref{sec:infer} 节，训练 infra 对应第~\ref{sec:train}、\ref{sec:cluster} 节，
算子/编译对应第~\ref{sec:kernel}、\ref{sec:compiler} 节，芯片岗对应第~\ref{sec:hw}、\ref{sec:accel} 节。

\begin{table*}[!t]
\centering
\caption{岗位 $\times$ 能力矩阵：工程与研究两条线的核心要求}
\label{tab:roles}
\footnotesize
\begin{tabularx}{\textwidth}{@{}L{2.0cm} Y Y L{3.6cm} L{3.4cm}@{}}
\toprule
岗位 & 工程能力（必须能做） & 研究能力（加分项） & 必须当场算出 & 常被深挖 \\
\midrule
LLM 推理系统 & KV 管理、batching 调度、量化部署、推测解码的工程取舍；读 profile 定位 decode 瓶颈 & 调度策略的理论分析；prefill/decode 分离的建模；新注意力变体的系统代价 & decode 步时、KV 容量、并发上限、TTFT/TPOT & vLLM 分页细节、chunked prefill、SLO 与 goodput \\
训练 infra & DDP/FSDP/3D 并行配置、显存预算、通信重叠、checkpoint 与故障恢复 & 并行策略自动搜索、通信压缩、故障与慢节点建模 & $16N$ 显存、all-reduce 时间、气泡、$6ND$ 天数 & ZeRO 各阶段、TP 为何限在节点内、MFU 排查 \\
算子 / 编译器 & CUDA/Triton kernel、tiling、融合、profile（Nsight） & 调度空间搜索、代价模型、动态形状 & 算术强度、tile 大小、HBM 字节、MFU & FlashAttention IO、bank conflict、occupancy \\
AI 芯片 / 加速器 & 数据流映射、片上存储规划、性能建模与仿真 & 新数据流与稀疏/低精度支持、软硬协同 & PE 利用率、片上带宽、能耗/字节 & systolic 填充开销、TPU 与 GPU 的取舍 \\
MLSys 研究 & 可复现实验、基线公平、消融设计 & 找到被忽视的前提、提出可证伪的假设 & 任一上界估算 & 论文的局限、你会做的下一个实验 \\
\bottomrule
\end{tabularx}
\end{table*}

\subsection{回答模板：结论 → 估算 → 瓶颈 → 权衡 → 证据}

开放式系统题最常见的失败是「列名词」：把知道的技术全部报一遍，却不说哪一个在当前场景下起作用。本文推荐的五步模板强制把回答落到数上：
\begin{enumerate}
\item \textbf{结论}：第一句话直接给判断（「这是带宽受限，先上 GQA 与 INT8 权重」）；
\item \textbf{估算}：用第~\ref{sec:perf} 节的 roofline 与附录的公式给出数量级；
\item \textbf{瓶颈}：说明估算显示的主导项是计算、访存还是通信，以及它会在什么条件下切换；
\item \textbf{权衡}：每个方案的代价（精度、时延、复杂度、显存）；
\item \textbf{证据}：如何测量来验证判断——指出能区分两个假设的那个观测。
\end{enumerate}

\begin{qa}{「把 Llama-3-70B 在线服务的成本降一半」，你怎么答？}
\ans \textbf{结论}：decode 受带宽限制，优先降低每 token 读的字节数，再提高 batch。\textbf{估算}：BF16 权重 141~GB，8 卡 TP 每卡每步读约 17.6~GB，
H100 3.35~TB/s 下约 5.3~ms/步；GQA 下每 token KV 320~KiB。\textbf{瓶颈}：batch 小时权重主导，batch 大、上下文长时 KV 主导。
\textbf{权衡}：INT8/FP8 权重把权重字节减半，精度损失需在任务上验证；FP8 KV 让并发翻倍；推测解码降低时延但大 batch 下收益递减；前缀缓存对系统提示共享的场景收益最大。
\textbf{证据}：先测当前 batch 分布与 KV 占用，若平均 batch 已经很大且 KV 占满，则 KV 量化与前缀缓存优先；若 batch 小，权重量化优先。
\wrong 不做估算直接列「量化、剪枝、蒸馏、FlashAttention、vLLM」。
\end{qa}

\subsection{常见错答}

表~\ref{tab:wrong} 汇总了前文面试追问框中出现频率最高的错答。它们几乎都源于同一个错误：把一个在特定区间成立的结论当成普遍规律，
却不说明区间的边界（batch、序列长度、拓扑、精度）。

\begin{table}[!t]
\centering
\caption{高频错答与纠正}
\label{tab:wrong}
\footnotesize
\begin{tabularx}{\columnwidth}{@{}L{3.0cm} Y@{}}
\toprule
错答 & 纠正 \\
\midrule
FlashAttention 把 attention 变成线性 & FLOPs 不变；HBM 流量 $\Theta(T^2d^2/M)$ 仍二次；显存从 $O(T^2)$ 到 $O(T)$ \\
decode 慢是因为算力不够 & 小 batch 下带宽受限；大 batch、短上下文才转为计算受限 \\
ZeRO-3 通信和 DP 一样 & 约为 DP 的 1.5 倍（多一次前向 all-gather） \\
ring all-reduce 时间随卡数线性增长 & 带宽项 $2(p-1)/p\cdot S/B$ 几乎与 $p$ 无关，延迟项线性 \\
7B 模型训练需要 14~GB & 混合精度 Adam 约 $16N=112$~GB，另加激活 \\
GQA 减少计算 & 主要减少 KV 字节，比例 $n_{\text{heads}}/n_{kv}$ \\
推测解码总能加速 & $\alpha$ 低或大 batch 时可能变慢 \\
MoE 省算力所以便宜 & 显存按总参数，且有 all-to-all 与负载不均 \\
利用率 100\% 说明 GPU 吃满 & 只表示有 kernel 在跑，看 SM 与 DRAM 吞吐 \\
\bottomrule
\end{tabularx}
\end{table}

\subsection{把配套实验变成 STAR 故事}

项目经历是面试中最可控的部分。本文配套仓库的每个实验都按「20 秒结论 → STAR → 白板 → 不能推出的结论 → 追问」组织，
表~\ref{tab:star} 列出了十个故事的结论与考察点。所有数字都来自 CPU 上的测量或解析/离散事件模型，不是 GPU 实测；
讲故事时必须主动说明这一点，并说出「这个实验\emph{不能}推出什么」——这恰恰是面试官判断你是否理解实验的最强信号。

\begin{table*}[!t]
\centering
\caption{本文配套实验的十个 STAR 故事（数字为 CPU 测量或模型输出）}
\label{tab:star}
\footnotesize
\begin{tabularx}{\textwidth}{@{}c L{2.4cm} Y L{3.4cm}@{}}
\toprule
\# & 实验 & 20 秒结论 & 考察点 \\
\midrule
1 & CS229S P1 FlashAttention & 按论文块大小，HBM 流量只比朴素少 2.00 倍（因果 3.7--4.0 倍），在 $T$ 上仍二次 & IO 复杂度、读论文 \\
2 & CS229S P2 推测解码 & 每次目标前向的 token 数与 $(1-\alpha^{\gamma+1})/(1-\alpha)$ 差 4\% 以内；CPU 上加速只有 0.98--1.99 倍 & 推理加速、口径 \\
3 & 15-849 服务 & 连续预留 KV 浪费 85.6\%，分页后 1.25\%，并发 104 → 714；chunked prefill 把最差一步 121~ms 压到 9.2~ms & serving \\
4 & 15-849 分布式 & ZeRO 0--3 与单进程差 $2.9\times10^{-16}$；ZeRO-3 通信恰为 DP 的 1.5 倍 & 分布式训练 \\
5 & 15-849 编译 & MLP epilogue 融合省 49\% HBM 流量，但只省 16\% 时间 & 编译器、roofline \\
6 & 15-849 压缩 & 8/1024 个 50 倍离群通道让 per-tensor W8A8 误差 2.3\% → 10.2\%，LLM.int8 拉回 0.49\% & 量化 \\
7 & CS149 A3 渲染 & 乱序画圆 93\% 像素错误；分块 + scan 保序后误差为 0 & 并行正确性 \\
8 & CS149 A1 SIMD & 每个 8 宽 gang 只要一条慢 lane，SIMD 利用率掉到 0.163 & 发散、负载均衡 \\
9 & CS149 A2 线程池 & 纯 spin 线程池在 CPython 中每次 launch 31--50~ms，sleep 版 38.7~µs & 同步、GIL \\
10 & CS149 A1 多进程 & 第一版测出 8 进程 21 倍「超线性」加速，根因是块大小改变了缓存行为 & 实验可信度 \\
\bottomrule
\end{tabularx}
\end{table*}

选择策略：只准备两个故事并讲透。推理岗用 1 与 3，训练岗用 4 与 6（或 4 与 10），算子/芯片岗用 1 与 7。先说结果数字，90 秒左右停下让对方打断。

\begin{qa}{讲一个你发现自己实验结论错了的经历。}
\ans （故事 10）第一版多进程 Mandelbrot 测出 8 个 worker 21 倍加速，在 4 物理核的机器上显然不可能。我没有直接报告，而是先怀疑基线：
块是按 worker 数切的，块变小后 numpy 的中间数组落进了缓存——同一张图串行整块算要 3.65~s，按小块串行算只要 1.29~s。
加速有一大半来自缓存行为改变，而不是并行。把所有调度的块大小统一后，加速回到与物理核数相符的范围。结论：比较两种实现时只能改变一个变量，超线性加速首先是测量问题的信号。
\wrong 只讲「我实现了什么」而没有「我怎么确认它是对的」。
\end{qa}

\begin{qa}{你的实验都在 CPU 上，怎么说明它对 GPU 有意义？}
\ans 区分三类结论：（1）正确性结论（ZeRO 逐位一致、分页 attention 与连续实现逐元素相等）与硬件无关；（2）计数结论（HBM 字节数、KV 浪费率、通信量）由算法决定，
与硬件无关；（3）时间结论只有相对关系与拐点位置有意义，绝对值需要在 GPU 上重测。我会主动说明哪个数属于哪一类，例如推测解码在 CPU 上加速只有 0.98--1.99 倍，
是因为 CPU 上验证 $\gamma+1$ 个 token 的代价是单 token 的 1.4--1.6 倍，而 GPU 上 batch 1 的 decode 受带宽限制，这个比值接近 1。
\wrong 把 CPU 上的墙钟加速直接当成 GPU 结论，或者反过来说「CPU 实验没有意义」。
\end{qa}

\subsection{为面试读论文}

读系统论文时，用五个问题做笔记，每个问题的答案不超过两句：
\begin{enumerate}
\item \textbf{瓶颈}：论文声称解决的瓶颈是计算、访存、通信还是调度/碎片？用式~\eqref{eq:three} 的哪一项表达？
\item \textbf{关键想法}：一句话说清它把什么换成了什么（例如 FlashAttention 用重算换 HBM 流量，vLLM 用间接寻址换连续性）。
\item \textbf{前提}：在什么区间成立？batch、序列长度、硬件代际、拓扑变了是否仍然成立？
\item \textbf{证据}：主结果的基线是否公平（同等调优、同一硬件）？加速来自哪一部分（是否有消融）？
\item \textbf{下一个实验}：如果你来做，会补哪个消融或扩展到哪个区间？
\end{enumerate}
以 vLLM~\cite{kwon2023vllm} 为例：瓶颈是 KV 的碎片（显存而非算力）；关键想法是把操作系统的分页搬到 KV 上；前提是请求长度方差大、显存是并发的约束；
证据是与 Orca 等系统的吞吐对比，2--4 倍的提升主要来自更大的 batch；下一个实验可以是在长度方差小（如固定长度批处理）的负载上，分页的收益是否消失。

\subsection{研究练习：提出一个消融}

研究岗面试常以「如果你来改进这篇论文，你会做什么实验」收尾。好的回答具备三个要素：一个可证伪的假设、一个只改变一个变量的实验、一个能区分假设的观测。

\begin{qa}{给 chunked prefill 设计一个消融实验。}
\ans \textbf{假设}：chunk 大小存在最优值——太小使 prefill 的 GEMM 不饱和、总吞吐下降，太大使与之同批的 decode 步变长、TPOT 尾延迟上升。
\textbf{实验}：固定负载（到达率、长度分布）与硬件，只扫 chunk 大小（128 到 8192），记录吞吐、TTFT 与 TPOT 的 P99 以及 goodput。
\textbf{区分性观测}：若 goodput 曲线呈倒 U 形且峰值位置随到达率移动，则支持假设并说明 chunk 大小应自适应；若单调，则说明当前负载下只有一侧的约束起作用。
本文配套的离散事件模拟器（第~\ref{sec:infer} 节）可以在 CPU 上先做这个扫描，再到 GPU 上验证峰值位置。
\wrong 「把 chunk 调大调小看看效果」——没有假设，也没有说明看哪个指标。
\end{qa}

\begin{qa}{给你一块没见过的新加速器，第一周你怎么评估它？}
\ans 先测三个上界：峰值算力（大方阵 GEMM）、内存带宽（stream 类拷贝）、互连带宽与延迟（不同消息大小的 all-reduce），得到它的 roofline 与 $\alpha$--$\beta$ 参数；
再用代表性负载（一个 decode 步、一个训练步）测实际达到上界的比例；最后看软件栈：编译器是否能自动融合、是否支持所需的精度与稀疏格式。
报告要写清：硬件、软件版本、预热与重复次数、中位数与离散度。
\wrong 直接跑一个端到端模型看速度，无法分辨慢在硬件、编译器还是负载。
\end{qa}

\subsection{实验与汇报的工程规范}

可信的数字有固定的产生方式：预热后多次测量取中位数并报告离散度；写明硬件、软件版本与线程数；测量前先用解析模型给出预期值，
偏差超过 20\% 时先怀疑测量再怀疑理论；比较两个实现时只改变一个变量；正确性先于性能（例如在 float64 下与参考实现比较到 $10^{-9}$）。
本文所有配套实验的数字都按此规范产生，原始 CSV 与生成脚本随仓库公开。

\begin{keypts}
\begin{itemize}
\item 回答模板：结论 → 估算 → 瓶颈 → 权衡 → 证据；第一句就给判断。
\item 错答多源于忽略适用区间；说清 batch、长度、拓扑、精度的边界。
\item 项目故事要带「不能推出什么」；CPU 实验的计数与正确性结论与硬件无关。
\item 读论文五问：瓶颈、关键想法、前提、证据、下一个实验。
\end{itemize}
\end{keypts}

\begin{research}
\begin{itemize}
\item 把本节的五步模板当作一个「可评分的协议」：对同一道题，模板化回答与自由回答在信息密度（每分钟给出的可验证数字）上差多少？这可以作为自我训练的度量。
\item 从表~\ref{tab:star} 任选一个实验，写出它在 GPU 上的复现计划：哪些数应该不变，哪些会变，预期变多少，为什么。
\end{itemize}
\end{research}
